
\subsubsection*{2.5.2}
\phantomsection
\label{sga3:VIIB:2.5.2}

Suppose that \(k\) is a perfect field. In this case, we saw in \SGARefLink{sga3:VIIB:1.6.2}{1.6.2}
that the morphism
\[
  p:X\longmapsto X_e
\]
has a section
\[
  s:X_e\longrightarrow X
\]
that depends functorially on \(X\). Consequently, when \(X\) is a
formal group, this section is a morphism of formal groups. We thus
obtain:

\medskip
\noindent\textbf{Proposition.}
\phantomsection
\label{sga3:VIIB:2.5.2:proposition}
{\itshape When \(k\) is a perfect field, every formal \(k\)-group
\(G\) is canonically isomorphic to the semidirect product of an
infinitesimal group \(G_0\) and an étale group \(G_e\) acting on
\(G_0\).}

If, moreover, \(G\) is commutative, then \(G\) is the product of
\(G_0\) and \(G_e\). By \SGARefLink{sga3:VIIB:2.2.2}{2.2.2}, this canonical decomposition of
commutative formal \(k\)-groups corresponds to an analogous
decomposition of affine commutative \(k\)-group schemes; see \SGARefLink{sga3:VIIB:2.5.3}{2.5.3}
below.

\medskip
\noindent\textbf{Remark 2.5.2.A.}
\phantomsection
\label{sga3:VIIB:2.5.2.A}
\footnote{\label{sga3:VIIB:nde:101}The editors added this remark.}
When \(k\) is not perfect, the exact sequence
\[
  1\longrightarrow G_0\longrightarrow G\longrightarrow G_e
  \longrightarrow 1
\]
need not split. Consider the following example, taken from
\([DG70]\), §III.6, 8.6. Let \(k\) be a nonperfect field of
characteristic \(p>0\), choose
\[
  \lambda\in k\setminus k^p,
\]
and let \(L_\lambda\) be the abelian restricted \(p\)-Lie algebra
with basis \((x,y)\), whose symbolic \(p\)-th power
\textup{(cf.\ VII\textsubscript{A}, \SGARefLink{sga3:VIIA:5.2}{5.2})} is given by
\[
  x^{(p)}=x,
  \qquad
  y^{(p)}=\lambda x.
\]
Let
\[
  U_p(L_\lambda)
  =
  k[x,y]/(x^p-x,\ y^p-\lambda x)
\]
be the restricted enveloping algebra of \(L_\lambda\)
\textup{(cf.\ VII\textsubscript{A}, \SGARefLink{sga3:VIIA:5.3}{5.3})}, and let \(G_\lambda\)
be the commutative formal \(k\)-group whose affine algebra is
\(U_p(L_\lambda)\).

Then \(G_\lambda\) is a nonsplit extension of the constant étale
\(k\)-group \((\mathbb Z/p\mathbb Z)_k\) by the infinitesimal
\(k\)-group
\[
  \alpha_{p,k}=\operatorname{Spf}k[x]/(x^p).
\]
\phantomsection
\label{sga3:VIIB:2.5.2.A:alpha-p-source-note}
\noindent\textit{Translator's note.}\enspace
The French re-edition prints
\(\operatorname{Spf}k[x]/(x^p-x)\) here. That algebra is étale,
whereas the asserted kernel is the standard infinitesimal group
\(\alpha_{p,k}=\operatorname{Spf}k[x]/(x^p)\). Indeed, the kernel
fiber \(x=0\) in \(U_p(L_\lambda)\) has coordinate algebra
\(k[y]/(y^p)\). The English text uses the standard infinitesimal
reading while preserving the printed reading in this note.
Thus there is a nonsplit exact sequence of commutative formal
\(k\)-groups:
\begin{center}
  \phantomsection
  \label{sga3:VIIB:2.5.2.A:nonsplit-formal-sequence}
  \(
0\longrightarrow\boldsymbol{\alpha}_{p,k}
\longrightarrow G_\lambda
\longrightarrow(\mathbb Z/p\mathbb Z)_k
\longrightarrow0.
\)

\end{center}

By duality, this corresponds to a nonsplit exact sequence of
commutative algebraic \(k\)-groups:
\begin{center}
  \phantomsection
  \label{sga3:VIIB:2.5.2.A:nonsplit-dual-sequence}
  \(
0\longrightarrow\boldsymbol{\mu}_{p,k}
\longrightarrow D'(G_\lambda)
\longrightarrow\boldsymbol{\alpha}_{p,k}
\longrightarrow0
\)

\end{center}
where
\[
  \mu_{p,k}=\operatorname{Spec}k[t]/(t^p-1).
\]
In this way one obtains all extensions of \(\alpha_{p,k}\) by
\(\mu_{p,k}\); see \([DG70]\), §III.6, 8.6.
